\paragraph{The regulated Lorentzian integral at finite \(N\) (conditional for \(N\ge3\)).}
Take traceless Hermitian \(N\times N\) matrices \(A_\mu\), \(\mu=0,\dots,9\), with indices contracted by \(\eta=\mathrm{diag}(-1,1,\dots,1)\), the quartic term \(S_{\mathrm b}=-\tfrac1{4g_{\mathrm{YM}}^2}\operatorname{Tr}[A_\mu,A_\nu][A^\mu,A^\nu]\) of Equation~\eqref{eq:Smatrix} and the fermion content of~\cite{IKKT1997}, an input since that equation is bosonic, with \(\operatorname{Pf}\mathcal M(A)\) the Pfaffian that the fermions give, \(\mathcal M\) being linear in \(A\). The oscillatory integral has been defined in three ways. (A) With a positive Gaussian factor~\cite{Steinacker2026Invitation}: \(I_\epsilon[O]=\int\mathrm dA\,O(A)\operatorname{Pf}\mathcal M(A)\,\mathrm e^{\mathrm iS_{\mathrm b}(A)-\epsilon\sum_\mu\operatorname{Tr}A_\mu^2}\), \(\epsilon>0\), over Hermitian matrices, which converges absolutely for polynomial \(O\). That reference has the opposite sign of the quartic term; since \(\operatorname{Pf}\mathcal M\) is real on Hermitian matrices, its integral is the complex conjugate of the one below for real \(O\). (B) As the limit \(u\to0^+\) on the contours \(\mathcal R_u\) below~\cite{Anagnostopoulos2026}. (C) By Faddeev--Popov fixing of the boosts with \(\operatorname{Tr}A_0A_i=0\)~\cite{AsanoNishimuraPiensukYamamori2025}. A Lorentz-invariant mass \(S_\gamma=-\tfrac{\gamma}{2g_{\mathrm{YM}}^2}\operatorname{Tr}A_\mu A^\mu\) may be added. The relation of these integrals to the Euclidean model is disputed. Reference~\cite{Steinacker2026Invitation} states that with the Gaussian factor the Lorentzian model is no longer related to the Euclidean model by Wick rotation; \cite{Anagnostopoulos2026} states equivalence for (B) at \(\gamma=0\) and its failure for \(\gamma>0\); and \cite{ChouNishimuraTripathi2025} reports inequivalence for (C) at \(N=2\) with four spatial matrices and no fermions.

\emph{Contours.} For \(0\le u\le1\) let \(\mathcal R_u\) be \(A_0=\mathrm e^{-3\pi\mathrm iu/8}\tilde A_0\), \(A_i=\mathrm e^{\pi\mathrm iu/8}\tilde A_i\), with \(\tilde A_\mu\) traceless Hermitian, and let \(S_{\mathrm E}(\tilde A)=\tfrac1{4g_{\mathrm{YM}}^2}\sum_{\mu,\nu}\|[\tilde A_\mu,\tilde A_\nu]\|^2\) in the Hilbert--Schmidt norm. On \(\mathcal R_u\) one has \(\operatorname{Re}\mathrm iS_{\mathrm b}=-\sin(\pi u/2)\,S_{\mathrm E}(\tilde A)\), and the regulator exponent has real part \(-\epsilon[\cos(3\pi u/4)\operatorname{Tr}\tilde A_0^2+\cos(\pi u/4)\sum_i\operatorname{Tr}\tilde A_i^2]\). For fixed \(\epsilon>0\) the integrand is holomorphic, and for \(u<2/3\) the boundary term of Stokes' theorem on a ball of radius \(R\) times \([0,u]\) is bounded by a polynomial in \(R\) times \(\mathrm e^{-\epsilon\cos(3\pi u/4)R^2}\). So \(I_\epsilon[O]\) equals the same integral on \(\mathcal R_u\) for every \(u<2/3\) (exact). On \(\mathcal R_1\) at fixed \(\epsilon\) the integrand is not absolutely integrable, since the regulator grows on \(\tilde A_0\) while \(S_{\mathrm E}\) vanishes on commuting matrices. At fixed \(\epsilon\) the contour therefore does not reach the Euclidean model.

\emph{The bound along the flat directions.} On \(\mathcal R_u\) the factor \(\operatorname{Pf}\mathcal M\) is a constant phase times the Pfaffian in \(\tilde A\) with the coefficient matrix of \(\tilde A_0\) multiplied by \(\mathrm e^{-\pi\mathrm iu/2}\). The convergence proof of~\cite{AustingWheater2001} is stated for the Euclidean Pfaffian, bounds its modulus, and uses of the coefficient matrices only the bound \(|\Gamma^\mu_{\alpha\beta}|\le1\) on their entries, as that reference says. The phase keeps that bound, and the step singling out one large matrix, written there for one index, becomes a finite sum over \(\mu\). The same proof therefore gives \(\int\mathrm d\tilde A\,(1+\|\tilde A\|)^k\,|\operatorname{Pf}\mathcal M(A(u,\tilde A))|\,\mathrm e^{-cS_{\mathrm E}(\tilde A)}<\infty\) for \(c>0\), \(0\le u\le1\) and \(k<14\). For \(N\ge3\) this extension rests on reading that proof, and no independent proof is given here. For \(N=2\) it is proved directly. With \(A_\mu=\sum_aA_\mu^a\sigma^a/2\) and \(\mathcal G_{ab}=\eta^{\mu\nu}A_\mu^aA_\nu^b\), one has \(\operatorname{Pf}\mathcal M=c_{\mathrm P}(\det\mathcal G)^4\), a constant \(c_{\mathrm P}\neq0\), for all complex \(A\): the Pfaffian is invariant under \(\mathrm{SO}(10,\mathbb C)\times\mathrm{SO}(3,\mathbb C)\), real Euclidean \(A\) has a singular-value form, the spinor branches into eight doublets of \(\mathrm{SO}(3)\), and the resulting \(6\times6\) determinant, of degree six, vanishes to second order in each singular value and is nonzero when all three equal \(1\) (also checked at random complex \(A\) to \(10^{-13}\)). By the Cauchy--Binet formula \(|\det\mathcal G(A(u,\tilde A))|\le\det\tilde{\mathcal G}\), with \(\tilde{\mathcal G}\) the Euclidean contraction of \(\tilde A\), and in the singular values the resulting integral converges exactly for \(k<14\).

\emph{Limit at finite \(N\).} Dominated convergence on \(\mathcal R_u\) at one fixed \(u\in(0,2/3)\), followed by Stokes' theorem from that \(u\) to \(u=1\) with the auxiliary factor \(\exp(-\delta[\mathrm e^{3\pi\mathrm i/8}\operatorname{Tr}A_0^2+\mathrm e^{-\pi\mathrm i/8}\sum_i\operatorname{Tr}A_i^2])\), whose exponent has negative real part on all these contours and which is removed as \(\delta\to0^+\) by the same domination, gives for polynomial \(O\) of degree \(k<14\)
\[
\lim_{\epsilon\to0^+}I_\epsilon[O]=\mathrm e^{7\pi\mathrm i(N^2-1)/4}\int\mathrm d\tilde A\;O\bigl(\mathrm e^{-3\pi\mathrm i/8}\tilde A_0,\mathrm e^{\pi\mathrm i/8}\tilde A_i\bigr)\operatorname{Pf}\mathcal M(-\mathrm i\tilde A_0,\tilde A_i)\,\mathrm e^{-S_{\mathrm E}(\tilde A)},
\]
with difference at most \(C_{O,N}\epsilon\) for \(k\le11\). Where the Euclidean partition function does not vanish, which is proved at \(N=2\) and is a hypothesis here for \(N\ge3\), ratios converge to Euclidean values times phases, for example \(\langle\operatorname{Tr}A_0^2\rangle\to\mathrm e^{-3\pi\mathrm i/4}\langle\operatorname{Tr}\tilde A_0^2\rangle_{\mathrm E}\) and \(\langle\sum_i\operatorname{Tr}A_i^2\rangle\to\mathrm e^{\pi\mathrm i/4}\langle\sum_i\operatorname{Tr}\tilde A_i^2\rangle_{\mathrm E}\), the relations stated for (B) in~\cite{Anagnostopoulos2026}. The same limit holds with \(\sum_{\mu,\nu}M_{\mu\nu}\operatorname{Tr}A_\mu A_\nu\) for any positive definite \(M\) in place of the regulator, since the argument then needs only \(u\) small. For these \(O\) the limit is therefore independent of the regulator and invariant under \(\mathrm{SO}(9,1)\) at finite \(N\), although each regulator breaks that group. It also holds for \(\gamma\to0\) with \(\gamma/\epsilon\) bounded above, the exponent then having real part \(-[\epsilon\cos(3\pi u/4)-\tfrac{\gamma}{2g_{\mathrm{YM}}^2}\sin(3\pi u/4)]\operatorname{Tr}\tilde A_0^2\) on \(\tilde A_0\) and the corresponding term with \(\pi u/4\) on \(\tilde A_i\). For fixed \(\gamma<0\) the contour argument reaches \(u=1\) with the mass retained. For fixed \(\gamma>0\) the admissible \(u\) closes like \(\epsilon/\gamma\) as \(\epsilon\to0\), and nothing is stated. At \(N=2\), reduced to three real variables (numerical), the values at fixed \(\epsilon\) and \(u=0.3,0.45,0.6\) agree to about eight digits, a direct evaluation on Hermitian matrices at \(u=0\) agrees within its grid error at the larger \(\epsilon\) tested, the Euclidean ratio of \(\sum_i\operatorname{Tr}\tilde A_i^2\) to \(\operatorname{Tr}\tilde A_0^2\) is \(9\), the limit is approached linearly in \(\epsilon\) with the slope given by the first moment, and as \(\epsilon\to0\) the moment \(\langle(\sum_i\operatorname{Tr}A_i^2)^{n}\rangle\) converges in modulus to the Euclidean value for \(2n\le12\), grows logarithmically for \(2n=14\) and grows as \(1/\epsilon\) for \(2n=16\).

\emph{Definition (C) at \(N=2\).} For at least four spatial matrices, write \(d\) for their number. The conditions \(\operatorname{Tr}A_0A_i=0\) force either \(A_0=0\), a sector which the Faddeev--Popov determinant removes as the delta functions are smeared and then sharpened, or \(\mathcal G=\mathrm{diag}(g_+,g_-,-g_3)\) up to rotation, with \(g_\pm,g_3\ge0\). The integrand is then \((g_+g_-g_3)^{(d-3)/2+p}|g_+-g_-|(g_3+g_+)(g_3+g_-)\) times the exponential, with \(p=4\) for the content above at \(d=9\) and \(p=0\) without fermions. The contour argument applies on the wedge \(g_+>g_-\), and at \(u=1\) it gives the Euclidean integrand with \(|g_+-g_3||g_--g_3|\) replaced by \((g_+-g_3)(g_--g_3)\). In an ordered triple the three positions of \(g_3\) give signs \(+,-,+\), which sum to \(1\) against the \(3\) of the Euclidean measure. So for \(\gamma\le0\), with a convergence factor as in (A) removed last, which (C) also needs since it is not absolutely convergent, and wherever the integrals converge and the partition function does not vanish, the expectation values of invariant observables of (C) at \(N=2\) equal those of the Euclidean model at the same phases (exact), namely the same \(O\) and \(S_\gamma\) on \(\mathcal R_1\) with weight \(\mathrm e^{-S_{\mathrm E}}\) and, for \(p=4\), the Pfaffian there; for \(p=0\), by the rotation \(A_\mu=\mathrm e^{\pi\mathrm i/8}\tilde A_\mu\), \(\mu=1,\dots,d+1\), with \(\tilde A_{d+1}\) the \(\tilde A_0\) above, this is the model with \(A_0=-\mathrm iA_{d+1}\) of~\cite{ChouNishimuraTripathi2025}. For \(d=4\), no fermions and \(\gamma<0\), this gives \(\langle\operatorname{Tr}A_\mu A^\mu\rangle\) of (C) equal to the Euclidean value, which grows as \(|\gamma|^{-1/2}\) when \(\gamma\to0^-\) in both models (numerical). The gap between the models that~\cite{ChouNishimuraTripathi2025} plots at \(\gamma<0\) is not reproduced by this reduction, and the difference is not resolved here. For \(\gamma>0\) no contour is available and that case is not treated.

For these observables, this order of limits, and \(\gamma\le0\) or \(\gamma\) as above, the regulated integral therefore carries at finite \(N\), conditionally for \(N\ge3\), no datum beyond the Euclidean model, with any mass term on \(\mathcal R_1\), and its phases. This is no statement about \(N\to\infty\) taken first, since the constants \(C_{O,N}\) are not shown to be uniform in \(N\); about fixed \(\gamma>0\); about degree \(14\) and above; or about non-polynomial observables such as \(\operatorname{Tr}\mathrm e^{\mathrm ik^\mu A_\mu}\), which grow on \(\mathcal R_u\), so that the domination above fails for them. It does not exclude a real-time reading obtained by continuation. It does not derive a signature, a background, a frame or the fermion content, and Assumption~\ref{ass:lorentzian} is not discharged. Problem~E remains open and PROGRAMME-CONDITIONAL.

